\documentclass[12pt,a4paper]{article}
\usepackage{amsmath}
\usepackage{amssymb}  % Provides \mathbb and other symbols
\usepackage{graphicx}
\usepackage{subcaption}
\usepackage{enumitem}
\usepackage{listings}
\usepackage{hyperref}
\usepackage{ulem}
\usepackage{xcolor}
\usepackage{amsthm}
\usepackage{tcolorbox}  % For colored boxes
\usepackage[table]{xcolor} % For coloring table cells
\usepackage{array}         % For better table alignment
\usepackage{mdframed}
\usepackage{cancel}
\usepackage[margin=2.5cm]{geometry}
\setlength{\parindent}{0pt}
\setlength{\parskip}{0.8em}

\begin{document}
\begin{titlepage}
\centering
\vspace*{\fill}
{\Huge Analysis III} \\[0.5cm]
{\Huge Home Work 4}\\[1.5cm]
{\Large Student: 26-712-224}\\[0.5cm]
{\Large \today}
\vspace*{\fill}
\end{titlepage}

\newpage
\textbf{Exercise 1 (Fubini's Theorem, 2+3+1 P.)}
Define $f:[0;1]^2\to\mathbb{R}$ by $f(0,0)=0$ and
\begin{align*}
f(x,y)=\frac{y^2-x^2}{(x^2+y^2)^2}
\end{align*}
for $(x,y)\neq(0,0)$.
\begin{enumerate}
\item[(a)] Prove that, for $(x,y)\in(0;1)^2$,
\begin{align*}
f(x,y)=\frac{\partial}{\partial x}\frac{x}{x^2+y^2}
=-\frac{\partial}{\partial y}\frac{y}{x^2+y^2}.
\end{align*}
\item[(b)] Compute the integrals
\begin{align*}
\int_0^1\left[\int_0^1 f(x,y)\,dx\right]dy
\qquad\text{and}\qquad
\int_0^1\left[\int_0^1 f(x,y)\,dy\right]dx.
\end{align*}
\item[(c)] Is $f\in L^1([0;1]^2;\mathbb{R})$? Justify.
\end{enumerate}

\vspace{1cm}
\noindent
Answer:\\
\\
(a) We have that:
\begin{align*}
\frac{\partial}{\partial x}x(x^2+y^2)^{-1}
&=(x^2+y^2)^{-1}
+x\left[\frac{-1}{(x^2+y^2)^2}\cdot2x\right]\\
&=\frac{1}{(x^2+y^2)^2}\left[x^2+y^2-2x^2\right]\\
&=\frac{y^2-x^2}{(x^2+y^2)^2}.
\end{align*}
Similarly:
\begin{align*}
\frac{\partial}{\partial y}y(x^2+y^2)^{-1}
&=\frac{x^2-y^2}{(x^2+y^2)^2}.
\end{align*}
i.e.
\begin{align*}
\frac{\partial}{\partial x}x(x^2+y^2)^{-1}
&=-\frac{\partial}{\partial y}y(x^2+y^2)^{-1} \\
&= f(x,y)
\end{align*}


\newpage
\textbf{Exercise 2 (Application of the MCT, 2+2 P.)}

Let $f \in L^1(\mathbb{R}^n;\mathbb{R})$ and denote by $B_\varepsilon(x)$ the ball of radius $\varepsilon > 0$ around $x \in \mathbb{R}^n$. We want to prove that
\begin{align*}
\lim_{\varepsilon \to 0^+} \int_{B_\varepsilon(x)} f(y)dy = 0.
\end{align*}
\begin{enumerate}
\item[(a)] Use the Monotone Convergence Theorem to prove the assertion under the additional assumption that $f \geq 0$.

\textit{Hint:} You may use without proof that $f\chi_B \in L^1(\mathbb{R}^n;\mathbb{R})$ for any ball $B \subset \mathbb{R}^n$. You may also assume that $\varepsilon$ decreases to $0$ monotonically.

\item[(b)] Prove the assertion without the additional assumption by reducing the general problem to the situation in (a).

\textit{Hint:} Sheet 3, Exercise 1 might be useful.
\end{enumerate}
\textit{Remark:} It is in general a good strategy to prove statements first for non-negative functions and then generalize.
\vspace{1cm}
\noindent
Answer:\\
\\

\end{document}